% ============================================================
% Secao 2: Arquitetura Geral
% ============================================================

\section{Arquitetura}
\label{sec:arquitetura}

O \aletheion{} segue a arquitetura \textit{decoder-only} transformer
\cite{vaswani2017attention, radford2019gpt2} com modifica\c{c}\~oes para
suportar o sistema epist\^emico intr\'inseco.

\subsection{Vis\~ao Geral}

Dado uma sequ\^encia de tokens $\bx = (x_1, \ldots, x_T)$, o modelo produz
duas sa\'idas simult\^aneas:

\begin{equation}
    \text{\aletheion{}}(\bx) = \big(\underbrace{\bm{l} \in \R^{B \times T \times V}}_{\text{logits}},\
    \underbrace{\mathcal{T} \in \R^{B \times T \times 30+}}_{\text{tomografia epist\^emica}}\big)
\end{equation}

onde $B$ \'e o \textit{batch size}, $T$ o comprimento da sequ\^encia, $V$ o
tamanho do vocabul\'ario, e $\mathcal{T}$ \'e a tomografia epist\^emica contendo
30+ campos por token.

\subsection{Pipeline de Forward Pass}

O forward pass segue 4 est\'agios:

\subsubsection{Est\'agio 1: Embeddings}

\begin{equation}
    \bh^{(0)} = \text{TokenEmb}(\bx) \cdot \sqrt{d} \in \R^{B \times T \times d}
\end{equation}

onde $d = d_{\text{model}}$ \'e a dimens\~ao do modelo. A sa\'ida do embedding
\'e escalada por $\sqrt{d}$ seguindo a pr\'atica padr\~ao de transformers.
O \aletheion{} utiliza
\textbf{RoPE} (\textit{Rotary Position Embedding}) \cite{su2021roformer}
aplicado dentro da aten\c{c}\~ao, em vez de embeddings posicionais aditivos:

\begin{equation}
    \text{RoPE}(\bm{q}, m) = \bm{q} \odot \cos(m\bm{\theta}) + \text{rotate}(\bm{q}) \odot \sin(m\bm{\theta})
\end{equation}

onde $m$ \'e a posi\c{c}\~ao e $\bm{\theta}_i = 10000^{-2i/d_{\text{head}}}$.

\subsubsection{Est\'agio 2: Transformer Blocks ($\times N$)}

Cada bloco segue o padr\~ao \textit{pre-norm} (estilo LLaMA/GPT-2):

\begin{align}
    \bm{a}^{(\ell)} &= \bh^{(\ell-1)} + \text{MultiHeadAttn}\big(\text{LN}(\bh^{(\ell-1)}), \text{RoPE}\big) \label{eq:attn} \\
    \bh^{(\ell)} &= \bm{a}^{(\ell)} + \text{FeedForward}\big(\text{LN}(\bm{a}^{(\ell)})\big) \label{eq:ffn}
\end{align}

\paragraph{Aten\c{c}\~ao Multi-Head com M\'ascara Causal.}
A aten\c{c}\~ao causal garante autorregress\~ao:

\begin{equation}
    \text{Attn}(\bm{Q}, \bm{K}, \bm{V}) = \text{softmax}\left(\frac{\bm{Q}\bm{K}^\top}{\sqrt{d_k}} + \bm{M}\right)\bm{V}
\end{equation}

onde $\bm{M}_{ij} = -\infty$ se $j > i$ (m\'ascara causal triangular inferior).

\paragraph{FeedForward com GELU.}
A camada feedforward utiliza a ativa\c{c}\~ao GELU \cite{hendrycks2016gelu}:

\begin{equation}
    \text{FFN}(\bx) = \bm{W}_2 \cdot \text{GELU}(\bm{W}_1 \bx)
\end{equation}

onde $\bm{W}_1 \in \R^{d \times d_{\text{ff}}}$ e
$\bm{W}_2 \in \R^{d_{\text{ff}} \times d}$, com $d_{\text{ff}} = 4d$.

\subsubsection{Est\'agio 3: Proje\c{c}\~ao de Sa\'ida}

\begin{equation}
    \bm{l} = \text{LN}_{\text{final}}(\bh^{(N)}) \cdot \bm{W}_{\text{emb}}^\top
\end{equation}

O \texttt{lm\_head} \'e \textbf{tied} com a matriz de embedding
($\bm{W}_{\text{emb}} \in \R^{V \times d}$), reduzindo par\^ametros e
melhorando generaliza\c{c}\~ao \cite{press2017tying}.

\subsubsection{Est\'agio 4: EpistemicHead}

Simultaneamente ao est\'agio 3, os \textit{hidden states}
$\bh^{(N)} \in \R^{B \times T \times d}$ e os padr\~oes de aten\c{c}\~ao
s\~ao passados ao \texttt{EpistemicHead}:

\begin{equation}
    \mathcal{T} = \text{EpistemicHead}\big(\bh^{(N)}, \bm{A}^{(1:N)}\big)
\end{equation}

Este \'e o componente central que diferencia o \aletheion{} de LLMs convencionais
e ser\'a detalhado na \Cref{sec:epistemico}.

\subsection{Hiperpar\^ametros do Modelo}

A \Cref{tab:config_base} apresenta a configura\c{c}\~ao do modelo de 354M
par\^ametros:

\begin{table}[H]
\centering
\caption{Configura\c{c}\~ao do modelo 354M.}
\label{tab:config_base}
\begin{tabular}{lrl}
\toprule
\textbf{Par\^ametro} & \textbf{Valor} & \textbf{Descri\c{c}\~ao} \\
\midrule
$d_{\text{model}}$ & 1024 & Dimens\~ao do modelo \\
$n_{\text{heads}}$ & 16 & N\'umero de heads de aten\c{c}\~ao \\
$d_{\text{head}}$ & 64 & Dimens\~ao por head ($d/n_h$) \\
$n_{\text{layers}}$ & 24 & N\'umero de blocos transformer \\
$d_{\text{ff}}$ & 4096 & Dimens\~ao feedforward \\
$T_{\max}$ & 1024 & Comprimento m\'aximo de sequ\^encia \\
$|V|$ & 50.257 & Tamanho do vocabul\'ario (tiktoken GPT-2) \\
Dropout & 0.1 & Taxa de dropout \\
\midrule
\multicolumn{3}{l}{\textit{Treinamento}} \\
Batch size & 4 & Micro-batch por GPU \\
Grad. accum. & 16 & Steps de acumula\c{c}\~ao \\
Effective batch & 64 & Sequ\^encias por update \\
Mixed precision & bf16 & Precis\~ao mista \\
Learning rate & $3 \times 10^{-4}$ & Com warmup cosine \\
\bottomrule
\end{tabular}
\end{table}

\subsection{Weight Tying}

O \aletheion{} utiliza \textit{weight tying} entre a camada de embedding e o
\texttt{lm\_head}:

\begin{equation}
    \bm{W}_{\text{lm\_head}} = \bm{W}_{\text{emb}}^\top
\end{equation}

Para um vocabul\'ario de 50.257 tokens com $d = 1024$, isso economiza
$\sim$51.5M par\^ametros (14.5\% do total).
